% s2-open.tex -- Section 2 of the reader's companion (drafter: aqft, R3; edited in R4).
% GJ3 facts appear only in the table's GJ3 column (context); each such cell is marked (R5).
% The "selection" cell of that column is a quotation of the record (main.tex l. 255),
% attributed to [R]; R5 checks it against main.tex, not against GJ3's pages.
\section{Why it was open}\label{sec:open}

\emph{Compactness gave Glimm and Jaffe a vacuum; it did not give convergence.}

Glimm and Jaffe constructed the model in the series I--IV \cite{GJ1,GJ2,GJ3,GJ4}. In
paper III, as the record recounts it, the cutoff states are averaged over space
translations and a $w^*$-compactness argument produces a limit point; whether the net
$(\omega_g)$ itself converges ``was not proved there'' \rec{\S1.1}{3}. The record quotes
the Expositions volume: the $g\to1$ limit of vacuum expectation values ``has not been
proved'' \cite[\S4]{GJExp}. It also cites Guerra, Rosen and Simon, who pose Euclidean
analogues (local $L^p$ bounds on the Euclidean vacuum density, local convergence of
Gibbs states) as Problems~1--2 \cite[pp.~125--126]{GRS75}; their note added on p.~128
records Newman's solution of the small-coupling $p=1$ versions, so ``the printed problems
must not be read as an unqualified current open-problem claim'' \rec{\S14}{32}.

Table~\ref{tab:gj3} sets the statements of GJ3 (context) beside those of the record.
What the record adds is in the rows \emph{limit} and \emph{selection}, and in the rate of
the row \emph{local norm}.

\begin{table}[ht]
\centering\small
\begin{tabular}{@{}>{\raggedright\arraybackslash}p{0.12\textwidth}
                >{\raggedright\arraybackslash}p{0.41\textwidth}
                >{\raggedright\arraybackslash}p{0.41\textwidth}@{}}
\toprule
property & GJ3 (context) & the record (weak coupling, \condtext{conditional on $\UGJS$})\\
\midrule
states &
% (R5)
$\omega_n$: averages of $\omega_{g_n}$ over space translations, $g_n=g(\cdot/n)$,
$g\equiv1$ on $[-3,3]$ \cite[(2.3)--(2.5), p.~210]{GJ3}
& $\omega_g$ for every $g\in\G$; no averaging \rec{Theorem~12.2}{30}\\[2pt]
limit &
% (R5)
a $w^*$ limit point; a subsequence of $(\omega_n)$ converges $w^*$ to it
\cite[Thm.~2.1, p.~211]{GJ3}
& the limit of the full net in each $\A(O)^*$ \rec{Theorem~12.2(1)}{30}\\[2pt]
local norm &
% (R5)
on each local algebra, $(\omega_n)$ lies in a norm-compact subset of the local dual
\cite[Thm.~2.3, p.~211]{GJ3}
& Cauchy in each $\A(O)^*$; rate $\Psi_{m_1}(d)$ when $g\equiv1$ on $O_d$
\rec{Theorem~11.1}{28}\\[2pt]
normality &
% (R5)
every limit point is normal on each local algebra \cite[Thm.~2.3, p.~211]{GJ3}
& $\omega_\infty|_{\A(O)}$ is normal \rec{Proposition~11.2}{29}\\[2pt]
energy &
% (R5)
in the representation of a $w^*$ limit point, the generator of time translations has
spectrum in $[0,\infty)$ \cite[Thm.~2.1, p.~211]{GJ3}
& $\omega_\infty$ is a ground state of $(\A,\alpha)$ in the spectral sense of
\rec{Proposition~12.1}{29}\\[2pt]
selection &
% (R5)
the record's reading: compactness ``selects a state but proves neither convergence nor
canonicity'' \rec{\S1.4}{4}
& every $w^*$ limit point of GJ3's states equals $\omega_\infty$ on each $\A(O)$
\rec{Corollary~12.3}{31}\\
\bottomrule
\end{tabular}
\caption{GJ3 and the record, statement by statement. The GJ3 column is context; its
last cell quotes the record.}
\label{tab:gj3}
\end{table}

Moments are not enough. Convergence of Schwinger functions at weak coupling follows
from the GJS cluster expansion \rec{\S1.1}{3}: along an affine path of Euclidean cutoffs
$h\in\CE$, GJS bound the change of the moments of products of smeared fields by
$O(1)e^{-m_1d/2}$, with $d$ the distance from the test functions to the change of cutoff,
eq.~(1.1.11) of \cite{GJS} as quoted in \rec{A8(vii$'$)(a)}{11}. In the record's words,
such convergence ``lives in the wrong topology'': ``there is no soft passage from moments
to norms on local von Neumann algebras'' \rec{\S1.4}{4}.

\begin{context}
The difference $V(g')-V(g)$ has infinite norm, so the machinery for gapped lattice
systems of \cite{Hastings,BMNS} ``does not start'' \rec{\S1.4}{4}.
\end{context}

Later work does not supply the norm statement either, as the record reads it
\rec{\S1.1}{3}. Albeverio--H\o egh-Krohn--Zegarlinski prove Euclidean Gibbs uniqueness
under their stated hypotheses \cite{AHZ}, and Bauerschmidt--Dagallier--Weber prove
uniqueness of the invariant measure for the $\Phi^4_2$ stochastic-quantization equation
above the critical temperature \cite{BDW}; neither is, by itself, a theorem about
local-norm convergence of this Hamiltonian cutoff net. The Hamiltonian renormalisation programme of
Rodriguez Zarate--Thiemann is restricted to finite volume \cite{RZT}. The record
identifies no theorem of local-norm convergence ``in the references reviewed here''; this
is not a priority claim.

\noindent\emph{In the record:} \S\S1.1 and 1.4, pp.~3--4; A8(vii$'$)(a), p.~11;
Theorem~11.1, Propositions~11.2 and 12.1, Theorem~12.2 and Corollary~12.3, pp.~28--31;
\S14, p.~32.
